% Generated by scripts/assemble.py; corrections belong in source/corrections.json.
\examhead{106}{中国科学技术大学 2026 年数学分析试卷}{104}
% p105-b24; source PDF page 105
\par\Needspace{5\baselineskip}
\blockanchor{p105-b24}
\majorquestion{一}计算题. 每题 10 分, 共 50 分.\par
% p105-b25; source PDF page 105
\blockanchor{p105-b25}
\question{1}求极限 $\lim_{n\to \infty}\sum_{k = 1}^{n}\frac{1}{n^2 + k}$\par
% p106-b00; source PDF page 106
\blockanchor{p106-b00}
\question{2}设 $x_0 \in \left(0, \frac{\pi}{2}\right)$, 且 $x_n = \sin x_{n-1} (n = 1, 2, \cdots)$, 求 $\lim_{n \to \infty} x_n, \lim_{n \to \infty} n x_n^2$.\par
% p106-b01; source PDF page 106
\blockanchor{p106-b01}
\question{3}设 $f(x) = \ln \frac{1 + x}{1 - x}$，求 $f^{(n)}(0)$.\par
% p106-b02; source PDF page 106
\blockanchor{p106-b02}
\question{4}设 $y = f(x)$ 是 $e^x + xy = e$ 在 (1,0) 点附近确定的隐函数, 求 $f''(0)$.\par
% p106-b03; source PDF page 106
\blockanchor{p106-b03}
\question{5}求 $z = \frac{1}{2} (x^2 + y^2)$ 被柱面 $(x^2 + y^2)^2 = 2xy$ 所截下的面积.\par
% p106-b04; source PDF page 106
\par\Needspace{5\baselineskip}
\blockanchor{p106-b04}
\majorquestion{二}（15分）设函数 $f(x,y) = \left\{ \begin{array}{ll}\frac{|xy|^{\alpha}}{x^{2} + y^{2}}, & (x,y)\neq (0,0);\\ 0, & (x,y) = (0,0). \end{array} \right.$\par
% p106-b05; source PDF page 106
\blockanchor{p106-b05}
(1) $\alpha$ 取何值时, $f(x,y)$ 在 $(0,0)$ 处连续?\par
% p106-b06; source PDF page 106
\blockanchor{p106-b06}
(2) $\alpha$ 取何值时, $f(x,y)$ 在 $(0,0)$ 处可微?\par
% p106-b07; source PDF page 106
\par\Needspace{5\baselineskip}
\blockanchor{p106-b07}
\majorquestion{三}(15 分) 设 $f(x)$ 在 $\mathbb{R}$ 上反常可积, 即 $\int_{-\infty}^{+\infty}|f(x)|\mathrm{d}x$ 收敛, 定义函数 $\varphi(x) = \int_{-\infty}^{+\infty}f(t)\cos (xt)\mathrm{d}t$, 证明: $\varphi(x)$ 在 $\mathbb{R}$ 上一致连续.\par
% p106-b08; source PDF page 106
\par\Needspace{5\baselineskip}
\blockanchor{p106-b08}
\majorquestion{四}(20 分) 设 $\Omega$ 为 $\mathbb{R}^3$ 中一有光滑边界的有界区域, 设 $u \in C^2(\Omega) \cap C^1(\overline{\Omega})$, 且 $\Delta u = 0$.\par
% p106-b09; source PDF page 106
\blockanchor{p106-b09}
(1) 证明: $\iint_{\partial \Omega} u \frac{\partial u}{\partial \pmb{n}} \, \mathrm{d}S = \iiint_{\Omega} \| \nabla u \|^{2} \, \mathrm{d}V.$\par
% p106-b10; source PDF page 106
\blockanchor{p106-b10}
(2) 若 $u|_{\partial \Omega} \equiv 1$, 证明: $u \equiv 1$.\par
% p106-b11; source PDF page 106
\par\Needspace{5\baselineskip}
\blockanchor{p106-b11}
\majorquestion{五}(15 分) 证明:\par
% p106-b12; source PDF page 106
\blockanchor{p106-b12}
(1) 对任意的 $x \in \mathbb{R}$, 有 $|\cos x| = \frac{2}{\pi} + \frac{4}{\pi} \sum_{n=1}^{\infty} \frac{(-1)^{n+1}}{4n^2 - 1} \cos(2nx)$.\par
% p106-b13; source PDF page 106
\blockanchor{p106-b13}
(2) $\lim_{\lambda \to +\infty}\int_0^1 f(x)|\cos \lambda x|\mathrm{d}x = \frac{2}{\pi}\int_0^1 f(x)\mathrm{d}x,$ 其中 $f(x)$ 在[0,1]上黎曼可积.\par
% p106-b14; source PDF page 106
\par\Needspace{5\baselineskip}
\par\Needspace{6\baselineskip}
\blockanchor{p106-b14}
\majorquestion{六}(15 分) 设 $f(x)$ 在 $[0, +\infty)$ 上可导, 且 $0 \leq f(x) \leq \frac{x}{1 + x^2}$, 证明: 存在 $\xi \in (0, +\infty)$, 满足\par
\nopagebreak[4]
% p106-b15; source PDF page 106
\blockanchor{p106-b15}
\displayformula{f ^ {\prime} (\xi) = \frac {1 - \xi^ {2}}{(1 + \xi^ {2}) ^ {2}}.}
% p106-b16; source PDF page 106
\par\Needspace{5\baselineskip}
\blockanchor{p106-b16}
\majorquestion{七}(10 分) 设 $f(x) \in C[0,1]$, 对任意的 $0 \leq a < b \leq 1$, $f$ 在 $[a,b]$ 上至少有两个不同的最大值点, 证明: $f(x)$ 为常值函数.\par
% p106-b17; source PDF page 106
\par\Needspace{5\baselineskip}
\par\Needspace{6\baselineskip}
\blockanchor{p106-b17}
\majorquestion{八}(10 分) 设 $f: [0,1] \to [0,1]$ 满足\par
\nopagebreak[4]
% p106-b18; source PDF page 106
\blockanchor{p106-b18}
\displayformula{| f (x) - f (y) | <   | x - y |, \forall x, y \in [ 0, 1 ], x \neq y.}
% p106-b19; source PDF page 106
\blockanchor{p106-b19}
证明：存在唯一的 $\xi \in [0,1]$，使得 $f(\xi) = \xi$\par
